\subsection{Algebraic Evolution Equations}\label{sec:aees}
For simple models, AEEs tend to be more computationally efficient
and conceptually easier to implement than models formulated in terms of ODEs
or DAEs. AEEs that aim to capture complex behaviors can quickly become
cumbersome to formulate and implement, but generally they do not incur
a truncation error like ODEs and DAEs and at their worst only may require an iterative root-finding procedure.
Because of the empirical nature of most AEEs, they tend to offer parameters that are more intuitive to practitioners and easier to identify from experimental data.
%
An AEE model is typically comprised of three components:
\begin{description}
    \item[Envelope] is specified as a functional relation between $\sigma$ and $\varepsilon$ which describes the monotonic behavior of the model.
    \item[Unloading] is specified as a set of rules which describe the hysteretic
    behavior of the model.
    \item[Degradation] is a rule for changing the shape of the envelope
    as energy is dissipated in the response.
\end{description}
%
An advantage of the AEE models is that they are generally parameterized by characteristic points on an envelope curve, which is associated with simple monotonic tension and compression tests.
Some standard envelopes are presented in the remainder of this section.

\subsubsection{Concrete}\label{sec:concrete}
Concrete envelope curves are formulated in terms of the following parameters (see also \cref{fig:confined-curve}):
\begin{table}[h]
\centering
\begin{tabular}{ll}
\toprule
\(f_{cp}\) & peak compressive stress \\
\(\varepsilon_{cp}\) & peak compressive strain \\
\(f_{cu}\) & compressive crushing strength \\
\(\varepsilon_{cu}\) & strain at crushing strength \\
\(\lambda\) & ratio between unloading slope at $\varepsilon_{cu}$ and
initial slope \\
\bottomrule
\end{tabular}
\end{table}

\begin{description}
\item[Hognestad]
  Hognestad's (\citet{hognestad1951study}, \citet{kent1973cyclic}, \textcite{scott1982stressstrain}) monotonic
  compression envelope is simply the piecewise union of an initially parabolic response, followed by a linear loss of strength: 
  \begin{equation}
  %\boxed{
  \sigma(\varepsilon)=\left\{\begin{aligned}
  K f_{cp}\left[2\left(\frac{\varepsilon}{\varepsilon_{cp}}\right)-\left(\frac{\varepsilon}{\varepsilon_{cp}}\right)^{2}\right], &\quad \varepsilon \leq \varepsilon_{cp} \\
  K f_{cp}\left[1-Z\left(\varepsilon-\varepsilon_{cp}\right)\right], &\quad \varepsilon_{cp}<\varepsilon \leq \varepsilon_{cu} \\
  0.2 K f_{cu}, &\quad \varepsilon > \varepsilon_{cu}
  \end{aligned}\right.
  %}
  \label{eq:hognestad}
  \end{equation}
  where $K$ (a scaling for strength) and $Z$ (strain softening slope) are model parameters. A hysteretic loading-reloading algorithm for this curve was proposed by \citet{yassin1994nonlinear}, which has been implemented in the widely used \texttt{Concrete02} material model of OpenSees.
%
\item[Popovics] The curve used by \cite{popov1978cyclic} is commonly known as the Mander curve and is given by the nonlinear expression: 
  \begin{equation}
%\boxed{
\begin{aligned}
\sigma_c(\varepsilon) &= f_{cp}\frac{(\varepsilon/\varepsilon_{cp})r}{r-1+(\varepsilon/\varepsilon_{cp})^r},
\end{aligned}
%}
  \label{eq:popovics}
\end{equation}
where $r$ is a parameter often computed from:
\begin{equation}
r =\frac{E_{c}}{E_{c}-\left(f_{c p} / \epsilon_{c p}\right)},
\label{eq:popovics1}
\end{equation}
where $E_c$ is the Young's modulus of the concrete material. 

% \item[Exponential] (see Concrete04)
% $$
% \begin{aligned}
% \sigma = f_{t}\beta^{\frac{\varepsilon      - e_{ct}}
%                           {\varepsilon_{tu} - e_{ct}}} \\
% e_{ct} = \frac{f_{t}}{E_{0}}
% \end{aligned}
% $$
\end{description}


\subsubsection{Steel}\label{sec:steel}
\begin{description}
    \item[Ramberg-Osgood] A popular envelope model for metalic materials was developed by \citep{ramberg1943description} who proposed the following equation:
    \begin{equation}
    \varepsilon=\frac{\sigma}{E}+a\left(\frac{\sigma}{\sigma_y}\right)^n,
    \label{eq:RO}
    \end{equation}
    where $E$ is the elastic modulus, $\sigma_y$ is the yield stress, assumed to be at the 0.2\% offset strain, $n$ is a parameter controlling the transition to yield, and $a$ is a yield offset parameter, typically taken as 0.002. Note that the stress \(\sigma\) is defined implicitly in terms of strain \(\varepsilon\) in this model and must be solved for numerically.
    \item[Goldberg-Richard] \cite{goldberg1963analysis} proposed a curve which furnishes the stress explicitly in terms of strain, as expressed below:
    \begin{equation}
    \bar{\sigma}(\bar{\varepsilon}) = \left[b{\bar{\varepsilon}} + \frac{(1-b){\bar{\varepsilon}}}{\left(1 + |{\bar{\varepsilon}}|^r\right)^\frac{1}{r}}\right],
    \label{eq:goldberg-richard}
    \end{equation}
    where $\bar{\sigma}=\sigma/\sigma_y$, $\bar{\varepsilon}=\varepsilon/\varepsilon_y$, $(\sigma_y, \varepsilon_y)$ is the yield point on the curve, $b$ is the strain hardening parameter, and the parameter $r$ influences the shape of the transition curve and takes account of the Bauschinger effect{\footnote{A property of materials where the material's stress/strain characteristics change as a result of the microscopic stress distribution of the material. For example, an increase in tensile yield strength occurs at the expense of the compressive yield strength in steel material.}}. A hysteretic loading-reloading algorithm for this curve was proposed by \cite{giuffre1970comportamento}, which was extended by \cite{filippou1983effects} to include isotropic hardening.
\end{description}

